\documentclass{article}
\usepackage[margin=1in]{geometry}
\usepackage{amssymb}
\usepackage{siunitx}
\usepackage{amsmath}
\usepackage{graphicx}
\usepackage{subcaption}
\usepackage{listings}
\usepackage{xcolor}
\usepackage{hyperref}
\usepackage[noabbrev,nameinlink]{cleveref}
\usepackage{booktabs}
\usepackage[american voltages]{circuitikz}

% Matplotlib's PGF backend uses this command for tick-label formatting.
\providecommand{\mathdefault}[1]{#1}

\sisetup{
  separate-uncertainty = true,
  per-mode = symbol
}

\lstset{
  basicstyle=\ttfamily\small,
  frame=single,
  columns=fullflexible,
  breaklines=true
}

\title{The Systems Behind Modern AI}
\author{
  Sam Ly \\
  Professor Thanh Nguyen \\
  CS 4630
}
\date{\today}

\begin{document}
\maketitle

\begin{abstract}
  Modern machine learning and AI runs primarily on Graphics Processing Units (GPUs),
  hardware initially designed to render graphics for video games. We briefly
  discuss the history of GPU architecture, and the justification behind some key innovations.
  AI models are composed of neural networks, which rely on the matrix multiplication as
  a computational primitive.
  Thus, efficient implementations of matrix multiplication are crucial in modern
  AI systems.
  To explore these implementations, we discuss the CUDA programming model,
  and develop the roofline model as a guiding principle for efficient hardware utilization.
  The demand for AI-specific compute has led to the development of numerous
  alternative parallel processing architectures such as Tensor Processing Units,
  Neural Processing Units, and Wafer Scale Engines.
  We conclude with a brief discussion of future evolutions of these systems and
  their justifications.
\end{abstract}

\end{document}